\section{The algebraic locus and the propagation criteria}\label{sec:propagation}

The results of \Cref{sec:basepoints} are pointwise: a class, a criterion and
a cycle at a point, and in \Cref{cor:exactform} an implication to be proved
at every point. This section studies the set of points at which the Weil class
is algebraic. Properness of the Hilbert scheme makes it a countable union of
closed subvarieties (\Cref{thm:closedlocus}), and rational isogenies make it
dense (\Cref{prop:densealg}); together these turn the propagation question
into a dichotomy (\Cref{cor:allornothing}). Two criteria for propagation
follow, a bound along one orbit (\Cref{thm:degreebound}) and a rank at one
point (\Cref{thm:semiregclosure}). We treat the bound first, together with
the transport of cycles along isogenies, which decides one of its two
readings, and with the limits of representatives built from products of
divisors; the rest of the section is devoted to the rank.

\subsection{The algebraic locus}
\label{ssec:properness}

\Cref{q:residual} is a question about the set
\[
  \Sigma \;=\; \bigl\{\,s\in\cD_{n,n} \;:\;
    \omega_{s}\ \text{is algebraic on}\ A_{s}\,\bigr\}.
\]
This set is restricted in shape, and the restriction changes the form
of the remaining problem: instead of a class to construct at every point, one
is left with a bound to prove along one orbit, or a rank to compute at one
point. The reason is that the Hilbert scheme is proper, so each piece of
$\Sigma$ cut out by bounded data is Zariski closed. The mechanism is standard:
the locus where a fixed Hodge class becomes algebraic is a countable union of
closed subvarieties, the companion of the statement that the Hodge locus
itself is one (see \cite[Chapter~5]{Voi03} and, in the form used in
\Cref{thm:cdk}, \cite{CDK95}). What is particular to the present situation is
that the class is determined by \Cref{thm:hdgcount} and written out by
\Cref{thm:explicitgens}, that its algebraic locus is nonempty in every family
by \Cref{thm:everyfamily}, and that the group acting on the family is known.

\begin{setupenv}[The quotient and the universal family]\label{setup:shimura}
Keep \Cref{setup:family}. Fix an integer $N\ge3$ and let $\Gamma$ be the
subgroup of $\SU(V,H)(\QQ)$ preserving the lattice $L\subset V$ used to build
the family and acting trivially on $L/NL$. Write
\[
  p\colon\cD_{n,n}\longrightarrow\cS=\Gamma\backslash\cD_{n,n} .
\]
Since $N\ge3$ the group $\Gamma$ is torsion-free, so it acts freely on
$\cD_{n,n}$, and by Baily and Borel \cite{BB66} the quotient $\cS$ is a smooth
quasi-projective complex variety; it is irreducible because $\cD_{n,n}$ is
connected. The level structure rigidifies, so $\cS$ carries a universal family
\[
  f\colon\cA\longrightarrow\cS,
\]
and the relative polarisation $3\eta$ is relatively very ample, which fixes a
relative projective embedding and hence a notion of degree in the fibres.
Write $G=\SU(V,H)$, an algebraic group over $\QQ$, so that
$G(\RR)$ acts transitively on $\cD_{n,n}$ and $\Gamma\subset G(\QQ)$.
\end{setupenv}

We record two facts that are used repeatedly below.

\begin{lemma}[The section descends]\label{lem:omegadescends}
The section $\omega$ of \Cref{setup:family} is $\Gamma$-invariant, hence
descends to a flat section of $R^{2n}f_{*}\QQ$ over $\cS$ that is of type
$(n,n)$ at every point.
\end{lemma}

\begin{proof}
The Weil line is $\HW=\bigwedge^{2n}_{K}V$, and an element of $\SU(V,H)$ is by
definition a $K$-linear automorphism of $V$ of $K$-determinant $1$. Its action
on the top $K$-exterior power is multiplication by that determinant, so
$\Gamma$ acts trivially on $\HW$ and in particular fixes $\omega$. Flatness is
then immediate, and the Hodge type at every point is
\Cref{prop:weilline}(ii).
\end{proof}

\begin{lemma}[One component is enough]\label{lem:onecomponent}
Let $(A,\iota,\eta)$ be of $(K,-1,n)$-Weil type with
$\iota(\cO_{K})\subseteq\End(A)$, and let $\HW\subset H^{2n}(A,\QQ)$ be its
Weil line. If one nonzero element of $\HW$ is algebraic, every element of
$\HW$ is algebraic.
\end{lemma}

\begin{proof}
Among the elements $\tau_{m}=m+\sqrt{-d}$ of $\cO_{K}$ with $1\le m\le7$ the
ratios $\tau_{m}/\bbar{\tau}_{m}$ are pairwise distinct, since
$\tau_{m}\bbar{\tau}_{m'}$ is rational only when $m=m'$. The group of roots of
unity in an imaginary quadratic field has order $2$, $4$ or $6$, so some $m$
in that range has $\tau=\tau_{m}$ with $\tau/\bbar\tau$ not a root of unity,
and then $\tau^{2n}\ne\bbar{\tau}^{2n}$, that is $\tau^{2n}\notin\QQ$, for
every $n\ge1$. A range is needed, since the witness $m=1$ fails for two
fields: for $d=1$ the ratio attached to $1+\sqrt{-1}$ is a primitive fourth
root of unity, and for $d=3$ the ratio attached to $1+\sqrt{-3}$ is a
primitive cube root of unity, so $(1+\sqrt{-d})^{2n}$ is rational when $d=1$
and $n$ is even, and when $d=3$ and $3\mid n$.

Since $\tau\in\cO_{K}$, the element $\iota(\tau)$ is an endomorphism of $A$,
so $\iota(\tau)^{*}$ on $H^{2n}(A,\QQ)$ is induced by the graph
correspondence and therefore carries algebraic classes to algebraic classes.
On $H^{1}$ it is multiplication by $\tau$, so on
$\bigwedge^{2n}_{K}V$ it is multiplication by $\tau^{2n}$. As
$\tau^{2n}\notin\QQ$ and $\HW$ is a one-dimensional $K$-vector space, the pair
$w,\;\tau^{2n}w$ is a $\QQ$-basis of $\HW$ for every $w\ne0$. Explicitly, fix
a $K$-basis $w_{0}$ of $\HW$ and write $\tau^{2n}=a+b\sqrt{-d}$, so that
$b\ne0$, and $w=(w_{1}+w_{2}\sqrt{-d})\,w_{0}$; the determinant of the pair
$w,\tau^{2n}w$ in the $\QQ$-basis $w_{0},\sqrt{-d}\,w_{0}$ is
$b(w_{1}^{2}+dw_{2}^{2})$, which is nonzero when $w\ne0$. If $w$ is
algebraic then so is $\iota(\tau)^{*}w=\tau^{2n}w$, and the two together span.
\end{proof}

\Cref{lem:onecomponent} writes out the last step of the proof of
\Cref{cor:whatisleft}, where the witness $\tau$ has to be chosen.

\begin{theorem}[The algebraic locus is a countable union of closed
subvarieties]\label{thm:closedlocus}
Keep \Cref{setup:shimura}. For a finite tuple $\underline{P}=(P_{1},\dots
,P_{m})$ of Hilbert polynomials of degree $n$ and a tuple
$\underline{q}\in\QQ^{m}$ put
\begin{align*}
  \cH_{\underline{P}}
  &\;=\;
  \Hilb^{P_{1}}(\cA/\cS)\times_{\cS}\cdots\times_{\cS}\Hilb^{P_{m}}(\cA/\cS),
  \\[2pt]
  \Sigma_{\underline{P},\underline{q}}
  &\;=\;
  \mathrm{pr}\Bigl(\bigl\{\,\underline{Z}\in\cH_{\underline{P}}
    \;:\;\textstyle\sum_{i}q_{i}[Z_{i}]=\omega\,\bigr\}\Bigr)
  \;\subseteq\;\cS .
\end{align*}
Then each $\Sigma_{\underline{P},\underline{q}}$ is Zariski closed in
$\cS$, and
\[
  p^{-1}\Bigl(\,\bigcup_{\underline{P},\underline{q}}
    \Sigma_{\underline{P},\underline{q}}\Bigr)\;=\;\Sigma ,
\]
the union running over a countable index set.
\end{theorem}

\begin{proof}
Each $\Hilb^{P_{i}}(\cA/\cS)$ is projective over $\cS$ by Grothendieck's
construction of the Hilbert scheme \cite{Gro61}, so the fibre product
$\cH_{\underline{P}}$ is
projective over $\cS$, and being of finite type it has finitely many connected
components.

Fix a connected component $T$, replace it by its reduction, which changes
neither the underlying space nor any cohomology group, and let
$\cZ_{i}\subset\cA\times_{\cS}T$ be the pullback of the $i$-th universal
subscheme. It is flat over $T$ and its fibres have Hilbert polynomial $P_{i}$
of degree $n$, so the associated $n$-cycles, the top-dimensional components
counted with their multiplicities, are the fibres of a cycle on
$\cA\times_{\cS}T$ whose class restricts on the fibre over $t$ to $[Z_{i,t}]$.
Restriction of a global class to the fibres of a topologically
locally trivial family is a flat section of the local system, so
$t\mapsto\sum_{i}q_{i}[Z_{i,t}]$ is flat on $T$; and $\omega$ is flat on $\cS$
by \Cref{lem:omegadescends}, hence flat after pullback. Two flat sections of a
local system over a connected base agree everywhere as soon as they agree at
one point, so the locus where they agree is a union of connected components of
$\cH_{\underline{P}}$. That locus is therefore itself projective over $\cS$,
and the image of a proper morphism is Zariski closed.

For the equality of sets, let $s\in\Sigma$. Then
$\omega_{s}=\sum_{i}q_{i}[Z_{i}]$ for finitely many closed integral subschemes
$Z_{i}\subset A_{s}$ of dimension $n$ and rationals $q_{i}$; taking
$P_{i}$ to be the Hilbert polynomial of $Z_{i}$ puts $p(s)$ in
$\Sigma_{\underline{P},\underline{q}}$. The converse holds because the class
of the $n$-cycle of a closed subscheme of dimension $n$ is algebraic. The
index set is countable, since a Hilbert polynomial is determined by its
finitely many integer coordinates in the binomial basis, and $\QQ$ is
countable.
\end{proof}

The bound on the Hilbert polynomials cannot be dropped: without it one takes
the image of an infinite disjoint union, and an infinite union of closed sets
need be neither closed nor constructible, and the dichotomy below is what
can be said of $\Sigma$ itself. Tuples are needed because a rational cycle is a rational combination
of several integral ones, which may have different Hilbert polynomials. The
relative Chow variety of effective $n$-cycles of bounded degree would serve
equally well, with the same proof, since only properness over the base and
local constancy of the cycle class are used.

\begin{corollary}[The dichotomy]\label{cor:allornothing}
Exactly one of the following holds.
\begin{enumerate}[label=\textup{(\alph*)},nosep]
\item $\Sigma=\cD_{n,n}$, and the Hodge conjecture holds in every degree for
  every abelian $2n$-fold of $(K,-1,n)$-Weil type with $\Hg(A)=\SU(V,H)$ in
  the family.
\item $\Sigma$ has empty interior in the analytic topology. Its complement is
  then a dense $G_{\delta}$, so $\omega_{s}$ fails to be algebraic for every
  $s$ outside a meagre set, and the Hodge conjecture fails.
\end{enumerate}
In particular, if $\omega_{s}$ is algebraic for all $s$ in some nonempty open
subset of $\cD_{n,n}$, then it is algebraic for all $s$.
\end{corollary}

\begin{proof}
Suppose $\Sigma$ contains a nonempty open set; then so does its image in
$\cS$, and we may pick a nonempty connected open $U\subseteq\cS$ contained in
that image. By \Cref{thm:closedlocus} the countably many sets
$\Sigma_{\underline{P},\underline{q}}\cap U$ are closed in $U$ and cover it.
A connected complex manifold is a Baire space, so one of them has nonempty
interior in $U$. A Zariski closed subset of the irreducible variety $\cS$
either is all of $\cS$ or is a proper closed subvariety, and a proper closed
subvariety of a connected complex manifold has empty interior. Hence that
$\Sigma_{\underline{P},\underline{q}}$ is all of $\cS$, so $\Sigma=\cD_{n,n}$
and (a) holds by \Cref{thm:residual}(a). If $\Sigma$ has empty interior, its
complement is the intersection of the countably many dense open sets
$\cD_{n,n}\setminus p^{-1}(\Sigma_{\underline{P},\underline{q}})$, and Baire
again gives that this intersection is dense. At each of its points
$\omega_{s}$ is a Hodge class by \Cref{lem:omegadescends} and is not
algebraic, so the Hodge conjecture fails there.
\end{proof}

\Cref{fig:properness} draws the two cases.

\begin{figure}[htbp]
\centering
\includegraphics[width=0.88\textwidth]{fig_properness}
\caption{The dichotomy of \Cref{cor:allornothing}. Each stratum
$\Sigma_{\underline{P},\underline{q}}$ is Zariski closed by
\Cref{thm:closedlocus}, and there are countably many of them; the base point
$s_{0}$ of \Cref{thm:everyfamily} lies on one. By \Cref{prop:densealg} the
union of the strata is dense, so every neighbourhood of a point $s$ lying
off all of them meets infinitely many strata. The conjecture for the family
is the assertion that no such $s$ exists, and by \Cref{cor:allornothing} it
is enough that one stratum has nonempty interior.}
\label{fig:properness}
\end{figure}

\begin{lemma}[Invariance under rational isogenies]\label{lem:Ginvariance}
$\Sigma$ is stable under the action of $G(\QQ)$ on $\cD_{n,n}$.
\end{lemma}

\begin{proof}
A point of $\cD_{n,n}$ is a complex structure $J$ on $V_{\RR}$ compatible with
$H$, and $g\in G(\QQ)$ carries $J$ to $gJg^{-1}$. The map $g$ is then
$\QQ$-rational and $\CC$-linear from $(V_{\RR},J)$ to $(V_{\RR},gJg^{-1})$, so
after clearing denominators it induces an isogeny $A_{J}\to A_{gJ}$, whose
graph is an algebraic correspondence. Algebraicity of a rational cohomology
class is preserved by such a correspondence and by its transpose. Finally $g$
lies in $\SU(V,H)$, so by the computation in \Cref{lem:omegadescends} it fixes
the Weil line pointwise; the correspondence therefore carries $\omega_{J}$ to
a nonzero rational multiple of $\omega_{gJ}$, and $\omega_{gJ}$ is algebraic
when $\omega_{J}$ is.
\end{proof}

Two density statements should be distinguished. For a class that does not
remain of Hodge type over the whole family, the locus where it is of Hodge
type is a countable union of proper closed subvarieties, each algebraic by
\Cref{thm:cdk}, and density there concerns the distribution of special
points. For the class $\omega$ of \Cref{setup:family} the Hodge locus is the
whole family by \Cref{prop:weilline}(ii), so \Cref{thm:cdk} gives no
information, and the set that carries information is the algebraic locus. The
next proposition concerns that set, and its proof uses no Hodge theory: the
algebraic locus is a union of orbits of a group whose rational points are
dense in its real points.

\begin{proposition}[The algebraic locus is dense]\label{prop:densealg}
Let $n\ge2$. Then $\Sigma$ is dense in $\cD_{n,n}$ in the analytic topology,
and its image in $\cS$ is dense.
\end{proposition}

\begin{proof}
The hermitian form $H$ has rank $2n\ge4$ over $K$ and signature $(n,n)$, so it
is isotropic over $\RR$; over every $\QQ_{p}$ a hermitian form of rank at
least three over a quadratic extension is isotropic; and hermitian forms over
an imaginary quadratic field satisfy the Hasse principle
\cite[Ch.~10, \S 6]{Sch85}. So $H$ is isotropic over $\QQ$ and $G$ is
$\QQ$-isotropic. The group $G=\SU(V,H)$ is simply connected and almost
$\QQ$-simple, so by Platonov's theorem $G(\QQ)$ is generated by the rational
points of the unipotent radicals of its proper parabolic $\QQ$-subgroups
\cite[Thm.~7.6 and \S 7.2]{PR94}. For a unipotent group $U$ over $\QQ$ the
exponential identifies $U(\QQ)$ with a $\QQ$-vector space inside the
corresponding $\RR$-vector space $U(\RR)$, so $U(\QQ)$ is dense in $U(\RR)$.
The closure of $G(\QQ)$ in $G(\RR)$ therefore contains $U(\RR)$ for every such
$U$, and those generate the identity component $G(\RR)^{+}$, which acts
transitively on $\cD_{n,n}$. Hence every $G(\QQ)$-orbit in $\cD_{n,n}$ is
dense.

By \Cref{thm:everyfamily} the set $\Sigma$ is nonempty, and by
\Cref{lem:Ginvariance} it is $G(\QQ)$-stable, so it contains a dense orbit.
The projection $p$ is continuous and surjective, so the image is dense.
\end{proof}

In particular, if case (b) of \Cref{cor:allornothing} occurs, then $\Sigma$
and its complement are both dense, and the two sets are distinguished by
Baire category and not by position.

\begin{corollary}[A Baire category criterion]\label{cor:meagre}
Let $n\ge2$. The Hodge conjecture holds for every abelian $2n$-fold of
$(K,-1,n)$-Weil type with $\Hg(A)=\SU(V,H)$ in the family of
\Cref{setup:family} if and only if $\Sigma$ is non-meagre in $\cD_{n,n}$, and
if and only if $\Sigma$ has nonempty interior.
\end{corollary}

\begin{proof}
If $\Sigma=\cD_{n,n}$, both conditions hold, and so does the conjecture, by
\Cref{thm:residual}(a). Otherwise case (b) of \Cref{cor:allornothing} holds,
so $\Sigma$ has empty interior; it is then the countable union of the closed
sets $p^{-1}(\Sigma_{\underline{P},\underline{q}})$, each with empty interior,
and so it is meagre. The points $s$ with $\Hg(A_{s})=\SU(V,H)$ form the
complement of a countable union of proper closed analytic subsets of $\cD_{n,n}$, by
Weil's theorem \cite{Weil77} recalled before \Cref{thm:hdgcount}, so by Baire
they are not all contained in the meagre set $\Sigma$. At a point of this
kind outside $\Sigma$ the Hodge conjecture fails for $A_{s}$.
\end{proof}

\subsection{The first criterion: a bound along one orbit}

The proof of \Cref{prop:densealg} produces the dense algebraic locus by
transporting one cycle along rational isogenies, and an isogeny multiplies
degrees. The dense set so obtained is therefore stratified by unbounded data,
and this is what prevents \Cref{cor:allornothing} from applying to it. A bound
on the data along one orbit removes the difficulty.

\begin{theorem}[A bound along the orbit closes the family]\label{thm:degreebound}
Let $n\ge2$ and let $s_{0}\in\Sigma$ be one of the base points supplied by
\Cref{thm:everyfamily}. Suppose that there is a finite set $\mathcal{T}$ of pairs
$(\underline{P},\underline{q})$ such that for every $g\in G(\QQ)$ the class
$\omega_{g\cdot s_{0}}$ is represented by some cycle whose Hilbert data lies
in $\mathcal{T}$. Then $\Sigma=\cD_{n,n}$, and the Hodge conjecture holds in every
degree for every abelian $2n$-fold of $(K,-1,n)$-Weil type with
$\Hg(A)=\SU(V,H)$ in that family.
\end{theorem}

\begin{proof}
The orbit $G(\QQ)\cdot s_{0}$ is dense by the proof of \Cref{prop:densealg},
and by hypothesis it is the union of the finitely many sets
$G(\QQ)\cdot s_{0}\cap p^{-1}(\Sigma_{\underline{P},\underline{q}})$ with
$(\underline{P},\underline{q})\in\mathcal{T}$. The union of the finitely
many closed sets $p^{-1}(\Sigma_{\underline{P},\underline{q}})$,
$(\underline{P},\underline{q})\in\mathcal{T}$, therefore contains a dense set,
and so it is all of $\cD_{n,n}$; in particular $\Sigma=\cD_{n,n}$. In fact one
of these sets is already all of $\cS$: each
$\Sigma_{\underline{P},\underline{q}}$ is Zariski closed in the irreducible
variety $\cS$, hence either equal to $\cS$ or nowhere dense, and a finite
union of nowhere dense sets is nowhere dense. The Hodge conjecture for the
family follows by \Cref{thm:residual}(a).
\end{proof}

\begin{remark}[Choice of the representing cycle]\label{rem:chosencycle}
We call the hypothesis of \Cref{thm:degreebound} the bounded criterion
\textup{(P1)}; \Cref{rem:tworemain} lists it beside the rank criterion
\textup{(P2)} of \Cref{thm:semiregclosure}. It quantifies over representing
cycles at each point of the orbit and asks only that, at each point, one of
them have bounded complexity. A stronger hypothesis would fix the cycle,
taking at $g\cdot s_{0}$ the transport of the cycle at $s_{0}$ along the
isogeny of \Cref{lem:Ginvariance}. \Cref{thm:p1forces} below shows that this
stronger form is never satisfied by a base cycle with a component of finite
stabiliser, since the transported data is then unbounded. The proof above
uses only the weaker hypothesis, which is what the covering argument needs.
\end{remark}

The hypothesis of \Cref{thm:degreebound} can be sharpened by locating the
part of the Hilbert data that is at issue. The leading coefficient is not.

\begin{proposition}[The leading coefficient is constant on the orbit]
\label{prop:leadingconstant}
Let $s_{0}\in\Sigma$ and $g\in G(\QQ)$, and let $\omega_{s}$ denote the Weil
class at $s$ and $\eta$ the polarisation class. Then
\[
  \int_{A_{g\cdot s_{0}}}\omega_{g\cdot s_{0}}\cdot\frac{\eta^{n}}{n!}
  \;=\;\int_{A_{s_{0}}}\omega_{s_{0}}\cdot\frac{\eta^{n}}{n!},
\]
so the degree of the transported class with respect to the polarisation, which
is the leading coefficient of the Hilbert polynomial of any cycle
representing it, is the same at every point of the orbit.
\end{proposition}

\begin{proof}
By \Cref{lem:Ginvariance} the correspondence transporting the class is induced
by a $\QQ$-rational map $g\in\SU(V,H)$. It preserves the hermitian form and
hence the polarisation class, it fixes the Weil line pointwise by
\Cref{lem:omegadescends}, and its determinant over $\QQ$ is
$\Nm_{K/\QQ}(\det_{K}g)=1$, so it acts trivially on the top exterior power of
$V^{\vee}$, in which both integrals are evaluated. The two integrals are
therefore equal. The identification of the number with the leading
coefficient is Hirzebruch--Riemann--Roch on $A$, whose Todd class is $1$ (for
the embedding by $3\eta$ of \Cref{setup:shimura} the leading coefficient is
$3^{n}$ times the number); the remaining coefficients involve the structure
sheaf of the cycle and not only its class.
\end{proof}

\begin{remark}[The source of unboundedness]
\label{rem:whereunbounded}
Write a representing cycle at a point of the orbit, the transported cycle or
any other, as a difference $\xi^{+}-\xi^{-}$ of effective cycles. By
\Cref{prop:leadingconstant} the difference of their degrees is constant along
the orbit, and by the boundedness of effective cycles of bounded degree in
the fibres of the polarised family $f\colon\cA\to\cS$, if $\deg\xi^{+}$ were
bounded along the orbit then only finitely many tuples
$(\underline{P},\underline{q})$ could occur and \Cref{thm:degreebound} would
apply. So the hypothesis of that theorem is equivalent to a bound on the
effective part alone.

This bound is not automatic, and the mechanism that threatens it can be seen
directly. The transport at $g$ is by the isogeny obtained from $g$ by
clearing denominators; if $g$ has denominator $m$ on the reference lattice,
which has rank $4n$, the isogeny has degree $m^{4n}$ (\Cref{lem:transport}
below), and pushing a cycle forward along it multiplies its multiplicities
and degrees by powers of $m$ (\Cref{lem:multiplicity}). Elements of $G(\QQ)$ of bounded denominator form a discrete subset of
$G(\RR)$, so along a dense orbit the denominators, and with them the isogeny
degrees, are unbounded. What \Cref{thm:degreebound} asks is that the cycles
nevertheless simplify: that at each point of the orbit the class $\omega$,
whose degree has not changed, admits an effective representation no more
complicated than before. For the transported cycle itself the answer is
negative as soon as a component of the base cycle has finite stabiliser
(\Cref{thm:p1forces} below), and for a cycle chosen at each point the bound
is equivalent to the conclusion (\Cref{thm:p1equivalent}). The present remark
shows only that the question reduces from the whole Hilbert datum to the
degree of the effective part, and that the part of the datum determined by
cohomology is constant.
\end{remark}

\subsection{Transport along isogenies}

The mechanism of \Cref{rem:whereunbounded} can be made explicit, and doing
so settles one of the two readings of \textup{(P1)}: the reading in
which the cycle at $s_{0}$ is transported along the orbit rather than chosen
afresh at each point. Write $\Lambda$ for the reference lattice, of rank $4n$
over $\ZZ$, and $E$ for the alternating form of the polarisation, so that a
point $s$ of $\cD_{n,n}$ is a complex structure on $\Lambda\otimes\RR$ and
$A_{s}=(\Lambda\otimes\RR)/\Lambda$. Throughout, $\dim A_{s}=2n$, and a
component of a cycle representing $\omega_{s}$ has dimension $n$.

\begin{lemma}[The transport in coordinates]\label{lem:transport}
Let $g\in G(\QQ)$ and let $c=c(g)$ be the least positive integer with
$cg\Lambda\subseteq\Lambda$. Then $\varphi_{g}:=cg$ induces an isogeny
$A_{s_{0}}\to A_{g\cdot s_{0}}$ with
\begin{enumerate}[label=\textup{(\roman*)},nosep,leftmargin=2.6em]
\item $\deg\varphi_{g}=c^{4n}$;
\item $\varphi_{g}^{*}\eta_{g\cdot s_{0}}=c^{2}\eta_{s_{0}}$ and
  $\varphi_{g}^{*}\omega_{g\cdot s_{0}}=c^{2n}\omega_{s_{0}}$;
\item the transported cycle representing $\omega_{g\cdot s_{0}}$ is
  $c^{-2n}\,\varphi_{g*}(Z)$, where $Z$ represents $\omega_{s_{0}}$.
\end{enumerate}
Moreover $\sup_{g\in G(\QQ)}c(g)=\infty$.
\end{lemma}

\begin{proof}
(i) The index of $cg\Lambda$ in $\Lambda$ is $|\det_{\QQ}(cg)|
=c^{4n}|\det_{\QQ}g|$, and $\det_{\QQ}g=\Nm_{K/\QQ}(\det_{K}g)=1$ for
$g\in\SU(V,H)$.

(ii) $g$ preserves $E$, so $(cg)^{*}E=c^{2}E$, which is the first identity.
On $H^{2n}=\bigwedge^{2n}(\Lambda^{\vee}\otimes\QQ)$ the map
$\varphi_{g}^{*}$ is $c^{2n}\bigwedge^{2n}g^{\vee}$, and $g\in\SU(V,H)$ fixes
the Weil line pointwise by \Cref{lem:omegadescends}.

(iii) $\varphi_{g*}\varphi_{g}^{*}=\deg\varphi_{g}$, so
$\omega_{g\cdot s_{0}}=(\deg\varphi_{g})^{-1}\varphi_{g*}
\varphi_{g}^{*}\omega_{g\cdot s_{0}}
=c^{-4n}\,c^{2n}\varphi_{g*}\omega_{s_{0}}=c^{-2n}\varphi_{g*}\omega_{s_{0}}$
by (i) and (ii).

For the last sentence, suppose that $c(g)\le C$ for every $g\in G(\QQ)$.
Then $(C!)\,g\Lambda\subseteq\Lambda$ for every $g$, so $G(\QQ)$ lies in the
discrete subset $(C!)^{-1}\End_{\ZZ}(\Lambda)$ of $\End(\Lambda\otimes\RR)$.
But by the proof of \Cref{prop:densealg} the closure of $G(\QQ)$ in $G(\RR)$
contains $G(\RR)^{+}$, a Lie group of positive dimension, which is a
contradiction.
\end{proof}

\begin{lemma}[Push-forward along an isogeny]
\label{lem:multiplicity}
Let $\varphi\colon A\to A'$ be an isogeny of abelian varieties and $Z\subseteq A$
an integral closed subvariety. Put
$G(Z)=\{x\in A:Z+x=Z\}$, a closed algebraic subgroup. Then
\[
  \varphi_{*}[Z]=m\,[\varphi(Z)],\qquad
  m=\bigl|\ker\varphi\cap G(Z)\bigr| .
\]
If moreover $\dim Z=n$ and $\varphi^{*}\eta'=c^{2}\eta$, then
\[
  \deg_{\eta'}\varphi(Z)=\frac{c^{2n}}{m}\,\deg_{\eta}Z .
\]
\end{lemma}

\begin{proof}
The map $\varphi|_{Z}\colon Z\to\varphi(Z)$ is finite and surjective, and its
degree $m$ is the number of points of a general fibre. For $z\in Z$ the fibre
through $z$ is $Z\cap(z+\ker\varphi)$, which as a set of translations is
$\{t\in\ker\varphi:z+t\in Z\}$. For a fixed $t$ the locus
$\{z\in Z:z+t\in Z\}=Z\cap(Z-t)$ is all of $Z$ when $t\in G(Z)$ and a proper
closed subset otherwise; since $\ker\varphi$ is finite, a general $z$ avoids
the finitely many proper subsets, and there
$\{t\in\ker\varphi:z+t\in Z\}=\ker\varphi\cap G(Z)$. This is the first
identity, and $\varphi_{*}[Z]=m[\varphi(Z)]$ by the definition of the
push-forward of cycles.

For the degree, the projection formula gives
\[
  m\deg_{\eta'}\varphi(Z)=\int\varphi_{*}[Z]\cdot\frac{\eta'^{n}}{n!}
  =\int[Z]\cdot\frac{(\varphi^{*}\eta')^{n}}{n!}
  =c^{2n}\int[Z]\cdot\frac{\eta^{n}}{n!}=c^{2n}\deg_{\eta}Z .
\]
\end{proof}

For example, for multiplication by $c$ on an abelian variety of dimension $2n$,
an abelian subvariety $Z$ of dimension $n$ has $G(Z)=Z$, so $m=|Z[c]|=c^{2n}$
and $\deg\varphi(Z)=\deg Z$, while a $Z$ whose stabiliser meets $A[c]$
trivially has $m=1$ and $\deg\varphi(Z)=c^{2n}\deg Z$.

\begin{theorem}[The transported cycle]\label{thm:p1forces}
Let $s_{0}\in\Sigma$ and let $Z=\sum_{i}q_{i}Z_{i}$, with $q_{i}\ne0$ and the
$Z_{i}$ integral of dimension $n$, represent $\omega_{s_{0}}$.
\begin{enumerate}[label=\textup{(\roman*)},leftmargin=2.6em,itemsep=2pt]
\item For $g\in G(\QQ)$ the transported cycle is
  $\sum_{i}c(g)^{-2n}m_{i}(g)\,q_{i}\,[\varphi_{g}(Z_{i})]$ with
  $m_{i}(g)=|\ker\varphi_{g}\cap G(Z_{i})|$, and the $i$-th component has
  degree $c(g)^{2n}\deg_{\eta}(Z_{i})/m_{i}(g)$.
\item If some $Z_{i}$ has finite stabiliser $G(Z_{i})$, then
  $\deg\varphi_{g}(Z_{i})\ge c(g)^{2n}\deg_{\eta}(Z_{i})/|G(Z_{i})|$, which is
  unbounded on $G(\QQ)$. The Hilbert data of the transported cycle then takes
  infinitely many values, and the hypothesis of \Cref{thm:degreebound} in the
  form that fixes the transported cycle is not satisfied by $Z$.
\item Consequently, a base cycle for which the transported data is bounded
  must have every component of positive-dimensional stabiliser, that is, every
  $Z_{i}$ must be stable under translation by a positive-dimensional abelian
  subvariety of $A_{s_{0}}$.
\end{enumerate}
\end{theorem}

\begin{proof}
(i) is \Cref{lem:transport}(iii) and \Cref{lem:multiplicity}.

(ii) $\ker\varphi_{g}\cap G(Z_{i})\subseteq G(Z_{i})$, so
$m_{i}(g)\le|G(Z_{i})|$, a constant; the displayed bound is (i), and $c(g)$ is
unbounded by \Cref{lem:transport}. Distinct degrees of the components give
distinct Hilbert polynomials, so the data takes infinitely many values.

(iii) is the contrapositive of (ii), together with the fact that an algebraic
subgroup of positive dimension in an abelian variety has an abelian subvariety
of positive dimension as its identity component.
\end{proof}

\begin{remark}[The two readings of \textup{(P1)}]\label{rem:p1meaning}
Three points should be kept apart.

First, \textup{(P1)} has two readings, and \Cref{thm:p1forces} settles one of
them. Read with the cycle fixed at $s_{0}$ and transported, it fails for every
base cycle written down in this paper: those cycles are intersections of
divisors, and an intersection of $n$ divisors on an abelian $2n$-fold has
finite stabiliser unless the divisors are pulled back from a quotient, which
the ones here are not. Read with the cycle chosen afresh at each point of the
orbit, which is the form of the hypothesis of \Cref{thm:degreebound}, it is
not affected by \Cref{thm:p1forces}.

Second, the difference between the two readings is the content of
\Cref{lem:multiplicity}. The class $\omega_{g\cdot s_{0}}$ has constant degree
along the orbit, by \Cref{prop:leadingconstant}; what grows is the complexity
of the particular representative that the transport produces, because the
isogeny spreads each component out by the factor $c^{2n}/m_{i}$. When a
component of the base cycle has finite stabiliser, the transport does not
produce a representative of bounded complexity, whether or not one exists.

Third, an approach to \textup{(P1)} must either produce a base cycle supported
on subvarieties with positive-dimensional stabiliser, a strong requirement for
a Weil class, or abandon the transport and bound instead the minimal
complexity of a representative, which is a statement about the Chow variety of
the family and not about any one cycle. Neither is achieved here.
\end{remark}

The bound on the minimal complexity of a representative is the hypothesis of
\Cref{thm:degreebound}, and it is the weakest hypothesis under which the
covering argument runs. The next theorem shows that it is equivalent to the
conclusion. Write $\mathbf{W}(K,n,\delta)$ for the statement that every
abelian $2n$-fold of $(K,-1,n)$-Weil type with discriminant $\delta$
(\Cref{def:disc}) has algebraic Weil classes, $\delta$ being the discriminant
of the family of \Cref{setup:family}.

\begin{theorem}[Equivalence of the bounded criterion]\label{thm:p1equivalent}
Let $n\ge2$, let $s_{0}\in\Sigma$ and keep \Cref{setup:family}. The following
are equivalent.
\begin{enumerate}[label=\textup{(\alph*)},leftmargin=2.6em,itemsep=2pt]
\item There is a finite set $\mathcal{T}$ of pairs
  $(\underline{P},\underline{q})$ such that at every point of the orbit
  $G(\QQ)\cdot s_{0}$ the class $\omega$ is represented by a cycle whose
  Hilbert data lies in $\mathcal{T}$; that is, \textup{(P1)} holds at
  $s_{0}$.
\item There is a single pair $(\underline{P},\underline{q})$ such that at
  every point of $\cD_{n,n}$, not merely of the orbit, the class $\omega$ is
  represented by a cycle with that datum.
\item $\Sigma=\cD_{n,n}$.
\item $\mathbf{W}(K,n,\delta)$ holds, that is, the Weil classes are algebraic
  on every member of the family.
\end{enumerate}
In particular \Cref{thm:degreebound} is an equivalence, and \textup{(P1)} is a
reformulation of the conjecture for the family and not a reduction of it. The
same holds for every CM field, with \Cref{thm:cmpropagation}(iv) in place of
\Cref{thm:degreebound}.
\end{theorem}

\begin{proof}
(a) implies (c) is \Cref{thm:degreebound}. (b) implies (a) is immediate, the
orbit being contained in $\cD_{n,n}$ and $\mathcal{T}$ being the one pair.
(d) implies (c) trivially, and (c) implies (d) because the algebraicity of
$\omega_{s}$ makes the whole Weil line of $A_{s}$ algebraic, by the argument
of \Cref{cor:whatisleft}, written out in \Cref{lem:onecomponent}.

It remains to deduce (b) from (c). If $\Sigma=\cD_{n,n}$ then $\Sigma$ has
nonempty interior, and the proof of \Cref{cor:allornothing} shows that some
$\Sigma_{\underline{P},\underline{q}}$ is then all of $\cS$: the countably
many sets $\Sigma_{\underline{P},\underline{q}}\cap U$ are closed in a
nonempty connected open $U\subseteq\cS$ and cover it, so by Baire one of them
has nonempty interior in $U$, and a proper Zariski closed subvariety of the
irreducible variety $\cS$ has empty interior. By the definition of
$\Sigma_{\underline{P},\underline{q}}$ in \Cref{thm:closedlocus} as the image of a proper morphism, every $s\in\cS$
then carries a tuple $\underline{Z}$ of subschemes with Hilbert polynomials
$\underline{P}$ and $\sum_{i}q_{i}[Z_{i}]=\omega$, which is (b).

For a CM field the same argument applies without change, using
\Cref{thm:cmpropagation}(iii) in place of \Cref{thm:closedlocus} and
\Cref{cor:allornothing}.
\end{proof}

\begin{remark}[Status of the first criterion]\label{rem:p1collapse}
By \Cref{thm:p1equivalent}, \textup{(P1)} is a restatement of the conjecture
for the family and not a route to it. \Cref{thm:degreebound} is a valid
implication, and its covering argument is the only place in the paper where
density of the rational orbit is converted into a global statement; but the
implication has a converse, so verifying its hypothesis for a family is as
hard as proving the conclusion for that family. The boundedness is
nevertheless a concrete finiteness statement about the Chow variety of the
universal family, and a reformulation of $\mathbf{W}(K,n,\delta)$ that
mentions neither Hodge theory nor deformation theory.

The two propagation criteria are therefore not on the same footing. The rank
criterion \textup{(P2)} of \Cref{thm:semiregclosure} is a sufficient
condition with no evident converse: injectivity of one semiregularity map at
one point implies the conjecture for the whole family, while
$\Sigma=\cD_{n,n}$ carries no information about the deformation theory of any
particular perfect complex. In its pure form, with Chern character a Weil
class, it is never satisfied (\Cref{thm:p2false}); the polarised criterion
\textup{(P2$'$)} of \Cref{rem:p2prime}, in which the Chern character may also
carry powers of the polarisation, is the form studied in
\Cref{sec:polarised}.
\end{remark}

\subsection{Products of divisors}

The intersection numbers needed below are Pfaffians. Let $E$ be a
nondegenerate alternating form on a real vector space $V$ of dimension $2g$,
let $\eta$ be its class, with the orientation for which
$e=\int\eta^{g}>0$, and for an $E$-self-adjoint endomorphism $\phi$ write
$x_{\phi}(u,v)=E(\phi u,v)$, an alternating form.

\begin{lemma}[Pfaffians of self-adjoint classes]\label{lem:pfaffian}
For real $s$, $\int(s\eta+x_{\phi})^{g}=e\,p_{\phi}(s)$, where $p_{\phi}$ is
the monic polynomial of degree $g$ with $p_{\phi}^{2}=\det(s+\phi)$. In
particular:
\begin{enumerate}[label=\textup{(\roman*)},nosep]
\item $2g\int x_{\phi}\eta^{g-1}=e\,\mathrm{tr}\,\phi$, and if
  $\mathrm{tr}\,\phi=0$ then
  $2g(g-1)\int x_{\phi}^{2}\eta^{g-2}=-e\,\mathrm{tr}(\phi^{2})$;
\item if $\mathrm{tr}\,\phi=0$ and $\phi^{2}=\beta$ is a nonzero real
  scalar, then $g$ is even and $\int x_{\phi}^{g}=e\,(-\beta)^{g/2}$.
\end{enumerate}
\end{lemma}

\begin{proof}
In a basis with Gram matrix $G$ for $E$, the class $s\eta+x_{\phi}$ has Gram
matrix $(s+\phi)^{T}G$, of determinant $\det G\det(s+\phi)$. Its Pfaffian
divided by $\mathrm{Pf}(G)$ is therefore a polynomial in $s$, monic of degree
$g$, whose square is $\det(s+\phi)$; and $\int\omega^{g}/e$ is
$\mathrm{Pf}(\omega)/\mathrm{Pf}(G)$ for every alternating form $\omega$.
This is the identity. Write $p_{\phi}(s)=s^{g}+p_{1}s^{g-1}+p_{2}s^{g-2}+\cdots$
and compare with
$\det(s+\phi)=s^{2g}+\mathrm{tr}(\phi)\,s^{2g-1}
+\tfrac12\bigl((\mathrm{tr}\,\phi)^{2}-\mathrm{tr}(\phi^{2})\bigr)s^{2g-2}+\cdots$:
then $2p_{1}=\mathrm{tr}\,\phi$ and
$2p_{2}+p_{1}^{2}=\tfrac12\bigl((\mathrm{tr}\,\phi)^{2}-\mathrm{tr}(\phi^{2})\bigr)$,
while the binomial expansion of $\int(s\eta+x_{\phi})^{g}$ gives
$g\int x_{\phi}\eta^{g-1}=e\,p_{1}$ and
$\binom{g}{2}\int x_{\phi}^{2}\eta^{g-2}=e\,p_{2}$; this is (i). Under the
hypothesis of (ii) the eigenvalues of $\phi$ are $\pm\sqrt{\beta}$, each with
multiplicity $g$ because the trace vanishes, so
$\det(s+\phi)=(s^{2}-\beta)^{g}$; this is the square of a polynomial only
for $g$ even, and then $p_{\phi}(s)=(s^{2}-\beta)^{g/2}$ and
$\int x_{\phi}^{g}=e\,p_{\phi}(0)$.
\end{proof}

At the quaternionic points of \Cref{setup:quat} the mechanism of
\Cref{rem:whereunbounded} can be seen explicitly. Take $n=2$ and the integral
lattice
$\Lambda_{b}=\cO_{b}^{2}$, $\cO_{b}=\ZZ\langle 1,i,j,ij\rangle$, for
$b\in\ZZ_{>0}$. Put
\[
  g_{b}(x,y)=E_{b}(jx,y),\qquad
  \lambda_{b}(x,y)=-\tfrac{1}{d}\,g_{b}(ix,y),\qquad
  Z_{b}=g_{b}^{2}-d\lambda_{b}^{2},\qquad Z'_{b}=2g_{b}\lambda_{b},
\]
so that $Z_{b}+\sqrt{-d}\,Z'_{b}=(g_{b}+\sqrt{-d}\,\lambda_{b})^{2}$ and
$Z_{b},Z'_{b}$ span the Weil plane. Write $x\in\cO_{b}$ as $x=u+vj$ with
$u,v\in\ZZ[i]$, where $i^{2}=-d$ as in \Cref{setup:quat}, so that
$\ZZ[i]=\ZZ[\sqrt{-d}]$; the pair $(\Lambda_{b},i)$ is then the same $\ZZ[i]$-module
$\ZZ[i]^{4}$ for every $b$, and we use this identification throughout.

\begin{proposition}[Divisibility of the quaternionic Weil cycle]
\label{prop:quatdivisible}
With these data, for every $b\in\ZZ_{>0}$:
\begin{enumerate}[label=\textup{(\roman*)},nosep]
\item $E_{b}=E_{u}+b\,E_{v}$, $g_{b}=b\,g_{1}$ and $\lambda_{b}=b\,\lambda_{1}$,
  where $E_{u},E_{v},g_{1},\lambda_{1}$ are integral on $\ZZ[i]^{4}$ and do not
  depend on $b$; in particular $g_{1},\lambda_{1}\in\mathrm{NS}(A_{b})$;
\item $Z_{b}=b^{2}Z_{1}$ and $Z'_{b}=b^{2}Z'_{1}$ in $\mathrm{CH}^{2}(A_{b})$
  modulo algebraically trivial cycles, for every choice of line bundles
  representing $g_{b},\lambda_{b}$; in cohomology the equalities hold in
  $\wedge^{4}\Lambda^{*}$;
\item $\int\eta_{b}^{4}=b^{2}\int\eta_{1}^{4}$,
  $\int g_{b}^{2}\eta_{b}^{2}=-\tfrac{b}{3}\int\eta_{b}^{4}$ and
  $\int Z_{b}^{2}=\tfrac{4b^{2}}{3}\int\eta_{b}^{4}$, so that
  $(\int g_{b}^{2}\eta_{b}^{2})^{2}=\tfrac{1}{12}\int Z_{b}^{2}\int\eta_{b}^{4}$;
\item for $\alpha\in K^{\times}$ the element $\alpha j$ generates the same
  algebra at the same point, and replacing $j$ by $\alpha j$ multiplies the
  two ratios in \textup{(iii)} by $N(\alpha)$ and $N(\alpha)^{2}$.
\end{enumerate}
\end{proposition}

\begin{proof}
For $x=u+vj$ one has $\bbar{x}=\bbar{u}-vj$, and the relations $ju=\bbar{u}j$,
$j^{2}=b$ give
$\bbar{x}\,i\,y\equiv\bbar{u}\,i\,u'+b\,v\,i\,\bbar{v'}$ modulo $Kj$. Since
$\mathrm{trd}$ vanishes on $Kj$ and restricts to $\mathrm{tr}_{K/\QQ}$ on $K$,
\[
  E_{b}(x,y)=\textstyle\sum_{k}a_{k}\bigl(\mathrm{tr}(i\,\bbar{u_{k}}u'_{k})
  +b\,\mathrm{tr}(i\,v_{k}\bbar{v'_{k}})\bigr).
\]
From $jx=b\bbar{v}+\bbar{u}j$ one reads off
$g_{b}(x,y)=b\sum_{k}a_{k}\bigl(\mathrm{tr}(i\,v_{k}u'_{k})
+\mathrm{tr}(i\,\bbar{u_{k}}\bbar{v'_{k}})\bigr)$, which is $b\,g_{1}$, and
$\lambda_{b}=b\lambda_{1}$ follows. The traces take integral values on
$\ZZ[i]$, and $g_{1}=b^{-1}g_{b}$ and $\lambda_{1}=b^{-1}\lambda_{b}$ are
rational multiples of Hodge classes, so they are integral Hodge classes and
lie in $\mathrm{NS}(A_{b})$ by the Lefschetz theorem on $(1,1)$-classes; this
proves (i). Then
$g_{b}+\sqrt{-d}\,\lambda_{b}=b(g_{1}+\sqrt{-d}\,\lambda_{1})$ and (ii) holds
in cohomology. Choose $L,M$ with $c_{1}(L)=g_{1}$ and
$c_{1}(M)=\lambda_{1}$. The cycle $c_{1}(L^{b})^{2}-d\,c_{1}(M^{b})^{2}$ equals
$b^{2}(c_{1}(L)^{2}-d\,c_{1}(M)^{2})$ in $\mathrm{CH}^{2}(A_{b})$, and another
choice changes both sides by algebraically trivial cycles.

For (iii), $E_{u}$ and $E_{v}$ live on complementary halves $U,V'$ of
$\Lambda\otimes\QQ$ of rank four, so $\eta_{u}^{3}=\eta_{v}^{3}=0$ and
$\int\eta_{b}^{4}=6b^{2}\int\eta_{u}^{2}\eta_{v}^{2}$. The form $g_{b}$ is
$x_{\phi}$ for $\phi=j$, which is $E_{b}$-self-adjoint by
\Cref{thm:quatweil}(i), anticommutes with $i$, so has trace $0$, and has
$j^{2}=b$; hence \Cref{lem:pfaffian}(i), with $g=4$, gives
$\int g_{b}^{2}\eta_{b}^{2}=-\tfrac{1}{24}\,\mathrm{tr}(j^{2})\int\eta_{b}^{4}
=-\tfrac{8b}{24}\int\eta_{b}^{4}$. For real $a,c$ the class
$ag_{b}+c\lambda_{b}$ is $x_{\phi}$ for $\phi=j(a-\tfrac{c}{d}\,i)$, again
self-adjoint of trace $0$, and
$\phi^{2}=(a+\tfrac{c}{d}i)\,j^{2}\,(a-\tfrac{c}{d}i)=b\,(a^{2}+c^{2}/d)$.
So \Cref{lem:pfaffian}(ii) gives
$\int(ag_{b}+c\lambda_{b})^{4}=b^{2}(a^{2}+c^{2}/d)^{2}\int\eta_{b}^{4}$, and
comparing coefficients, $\int g_{b}^{4}=b^{2}\int\eta_{b}^{4}$,
$\int g_{b}^{2}\lambda_{b}^{2}=\tfrac{b^{2}}{3d}\int\eta_{b}^{4}$ and
$\int\lambda_{b}^{4}=\tfrac{b^{2}}{d^{2}}\int\eta_{b}^{4}$. Hence
$\int Z_{b}^{2}=\int g_{b}^{4}-2d\int g_{b}^{2}\lambda_{b}^{2}
+d^{2}\int\lambda_{b}^{4}=\tfrac{4b^{2}}{3}\int\eta_{b}^{4}$.

For (iv), $(\alpha j)^{2}=N(\alpha)b$, and
$g_{\alpha j}+\sqrt{-d}\,\lambda_{\alpha j}$ is $\alpha$ or $\bbar{\alpha}$
times $g_{j}+\sqrt{-d}\,\lambda_{j}$.
\end{proof}

\begin{remark}[Growth of the polarisation]\label{rem:quatdivisible}
At these points no multiple growing with $b$ enters when the Weil class is
written in integral terms: by \Cref{prop:quatdivisible}(ii) the cycle $Z_{b}$
is $b^{2}Z_{1}$, with $Z_{1}$ one integral algebraic cycle built from the two
fixed divisor classes $g_{1},\lambda_{1}$. The factor that the orbit argument
would otherwise have to absorb into the coefficients $\underline{q}$ of
\Cref{thm:degreebound} is therefore absent. What grows is the polarisation:
$\int\eta_{b}^{4}=b^{2}\int\eta_{1}^{4}$, so the points $A_{b}$ carry
polarisations of unbounded degree, and an isogeny carrying $\eta_{b}$ to the
polarisation of one fixed family has degree $\int\eta_{b}^{4}/\int\eta^{4}$,
which grows like $b^{2}$. This is the unbounded isogeny degree of
\Cref{rem:whereunbounded}, located at an explicit family of points, and the
integrality of $Z_{1}$ does not affect it.
\end{remark}

The base points of \Cref{thm:everyfamily}, and every algebraicity statement
for Weil classes proved in this paper, represent $\omega$ through products of
divisor classes. The next result shows that a representation of that kind
can hold with bounded degrees only on a nowhere dense set, so the bound asked
for in \Cref{thm:degreebound} has to come from cycles that are not products of
divisors. For a class $x\in H^{2}(A_{s},\RR)$ put
\[
  t(x)=\frac{\int x\,\eta^{2n-1}}{\int\eta^{2n}},\qquad
  I(x)=t(x)^{2}\!\int\eta^{2n}-\int\bigl(x-t(x)\eta\bigr)^{2}\eta^{2n-2}.
\]
Since $\cD_{n,n}$ is contractible, the local system $R^{2}\pi_{*}\ZZ$ is
trivial; write $H^{2}_{\ZZ}$ for its fibre, identified with $H^{2}(A_{s},\ZZ)$
at every $s$. Both $t$ and $I$ are polynomial functions on $H^{2}_{\ZZ}$ and
do not depend on $s$.

\begin{lemma}[The Hodge norm of a divisor]\label{lem:nefnorm}
Let $h_{s}$ be the Hodge norm on $H^{2}(A_{s},\RR)$ defined by the
polarisation, normalised by $h_{s}(\eta)=\int\eta^{2n}$. Then
\begin{enumerate}[label=\textup{(\roman*)},nosep]
\item $h_{s}(x)=I(x)$ for every real class $x$ of type $(1,1)$ at $s$;
\item if $D$ is an effective divisor on $A_{s}$ then
  $I([D])\le 2\bigl(\int[D]\eta^{2n-1}\bigr)^{2}\big/\!\int\eta^{2n}$.
\end{enumerate}
\end{lemma}

\begin{proof}
Write $x=t\eta+x_{0}$ with $t=t(x)$, so that $x_{0}$ is primitive. The
Lefschetz decomposition is orthogonal for $h_{s}$, and on primitive classes
$h_{s}(x_{0})=-\int x_{0}\wedge C_{s}x_{0}\wedge\eta^{2n-2}$, where $C_{s}$ is
the Weil operator, positive by the Hodge--Riemann relations. For $x_{0}$ of type
$(1,1)$, $C_{s}x_{0}=x_{0}$, which gives (i). An effective divisor on an
abelian variety is nef, since a translate of it meets any given curve
properly, so $\int[D]^{2}\eta^{2n-2}\ge0$. Expanding,
$\int[D]^{2}\eta^{2n-2}=t^{2}\int\eta^{2n}+\int x_{0}^{2}\eta^{2n-2}$, hence
$-\int x_{0}^{2}\eta^{2n-2}\le t^{2}\int\eta^{2n}$ and
$I([D])\le2t^{2}\int\eta^{2n}$, which is (ii).
\end{proof}

\begin{proposition}[Divisor products of bounded degree]\label{prop:divroute}
In \Cref{setup:family}, fix $C>0$ and let $\Sigma_{C}\subset\cD_{n,n}$ be the
set of points $s$ at which $\omega_{s}$ is a $\QQ$-linear combination of
classes $[D_{1}]\cdots[D_{n}]$ with every $D_{i}$ an effective divisor on
$A_{s}$ of degree $\int[D_{i}]\eta^{2n-1}\le C$. Then $\Sigma_{C}$ is
contained in a locally finite union of proper closed analytic subsets of
$\cD_{n,n}$. In particular $\Sigma_{C}$ is nowhere dense, and on any dense
subset of $\cD_{n,n}$ at which $\omega$ is a combination of products of
effective divisors the degrees of those divisors are unbounded on every
nonempty open set.
\end{proposition}

\begin{proof}
For $x\in H^{2}_{\ZZ}$ let $\mathrm{Hdg}(x)\subset\cD_{n,n}$ be the
closed analytic set where $x$ is of type $(1,1)$. If $x\notin\QQ\eta$ it is
proper, because $\NS(A_{s})\otimes\QQ=\QQ\eta$ at a very general point by
\Cref{lem:genericNS}. Let $s\in\Sigma_{C}$. Not every $[D_{i}]$ in the
representation lies in $\QQ\eta$, since $\QQ\eta^{n}$ and the Weil line carry
different characters by \Cref{lem:charsdiffer}. So $s\in\mathrm{Hdg}(x)$ for
some $x$ in
\[
  X_{C}=\bigl\{x\in H^{2}_{\ZZ}\setminus\QQ\eta\;:\;
  I(x)\le 2C^{2}/\textstyle\int\eta^{2n}\bigr\},
\]
by \Cref{lem:nefnorm}(ii). Now let $U\subset\cD_{n,n}$ be open with compact
closure and $s_{0}\in U$. The Hodge norm depends continuously on the point, so
$h_{s_{0}}\le c_{U}\,h_{s}$ for some $c_{U}$ and all $s\in U$. If
$\mathrm{Hdg}(x)$ meets $U$ at $s$ and $x\in X_{C}$, then by
\Cref{lem:nefnorm}(i)
$h_{s_{0}}(x)\le c_{U}h_{s}(x)=c_{U}I(x)\le 2c_{U}C^{2}/\int\eta^{2n}$, and
only finitely many integral classes satisfy this. So the sets
$\mathrm{Hdg}(x)$, $x\in X_{C}$, form a locally finite family of proper closed
analytic subsets. Their union is closed with empty interior in the connected
manifold $\cD_{n,n}$, and it contains $\Sigma_{C}$.
\end{proof}

\begin{remark}[Divisorial representatives]\label{rem:divroute}
The loci on which this paper proves $\omega$ algebraic are the split,
quaternionic and product loci and their translates under $G(\QQ)$, and on
each of them $\omega$ is a rational polynomial in divisor classes, since an
isogeny carries $\NS\otimes\QQ$ isomorphically. \Cref{prop:divroute} shows
that no bounded family of such cycles covers a dense set, so the finiteness
in \Cref{thm:degreebound} cannot be obtained by bounding the divisors that
produce $\omega$ at the points where it is known to be algebraic. It would
have to come from representatives of another kind, such as the transported
cycles of \Cref{lem:Ginvariance}, which are push-forwards along isogenies
rather than products of divisors on the target; and by \Cref{thm:p1forces}
their data is unbounded as soon as a component of the base cycle has finite
stabiliser. The unbounded isogeny degree of \Cref{rem:whereunbounded} and
\Cref{rem:quatdivisible} is thus unavoidable for divisorial constructions.
The argument uses nothing about $\omega$ beyond the fact that it is not a
multiple of $\eta^{n}$, and nothing about the family beyond the generic
Picard number, so it applies equally to any variational argument that rests
on products of divisors.
\end{remark}

At $n=2$ the phenomenon can be exhibited on one fixed polarised lattice. Take
$\Lambda=\cO^{2}$ with $\cO=\ZZ\langle1,i,j,ij\rangle$, $i^{2}=-d$, $j^{2}=1$,
and $E$ as in \eqref{eq:quatE} with $a_{1}=a_{2}=1$. The lattice and the
polarisation are fixed, and all the loci below lie in the one family
$\cD_{2,2}$ that they define. Let
$R=\{x:\ x(iu,v)=x(u,iv)\}\subset H^{2}(\Lambda,\QQ)$ be the space of
$K$-bilinear classes, where the divisor classes of the quaternionic points
live, and for $x\in R$ let $\phi_{x}=E^{-1}x$, so that $x(u,v)=E(\phi_{x}u,v)$.
Put $e=\int\eta^{4}$ and, for $k\in\ZZ$,
\[
  C_{k}=\begin{pmatrix}1&k\\k&-1\end{pmatrix},\qquad
  \phi_{k}(u_{1},u_{2})=\tfrac{1}{2d}\bigl(j(u_{1}+ku_{2}),\,j(ku_{1}-u_{2})\bigr),
  \qquad x_{k}(u,v)=E(\phi_{k}u,v),
\]
so that $\phi_{k}$ is left multiplication by the matrix $(2d)^{-1}C_{k}\,j$
and $x_{k}=x_{0}+k(x_{1}-x_{0})$ is a pencil.

\begin{proposition}[A pencil of quaternionic loci]\label{prop:pencil}
In this setting, for every $d$ and every $k\in\ZZ$:
\begin{enumerate}[label=\textup{(\roman*)},nosep]
\item $e=384\,d^{4}$; $R$ has dimension $12$, every $x\in R$ has $t(x)=0$
  and $I(x)=\tfrac{1}{24}\,\mathrm{tr}(\phi_{x}^{2})\,e$, and the form $I$ has
  signature $(8,4)$ on $R$;
\item $x_{k}$ is a primitive integral class in $R$, and
  $\phi_{x_{k}}^{2}=\beta(k)$ with $\beta(k)=(1+k^{2})/(4d^{2})$;
\item the set $\mathrm{Hdg}(x_{k})$ is a quaternionic locus with algebra
  $(-d,\beta(k))\cong(-d,1+k^{2})$, nonempty of dimension $3$, and loci with
  different values of $\beta$ are not related by any automorphism of
  $(\Lambda,E,i)$;
\item at a very general point of $\mathrm{Hdg}(x_{k})$ one has
  $\NS=\ZZ\,\tfrac{1}{2d}\eta\oplus\ZZ x_{k}\oplus\ZZ\,\iota x_{k}$, where
  $\iota x=x(i\cdot,\cdot)$, and the minimum of $I$ on $\NS\setminus\QQ\eta$
  is $\mu_{k}=32\,d^{2}(1+k^{2})$.
\end{enumerate}
Consequently, at a very general point of $\mathrm{Hdg}(x_{k})$ every effective
divisor $D$ whose class is not a multiple of $\eta$ has
$\int[D]\eta^{3}\ge32\sqrt{6}\,d^{3}\sqrt{1+k^{2}}$, which grows linearly in
$|k|$.
\end{proposition}

\begin{proof}
Two facts about $\cO$ are used throughout. The reduced trace form
$\mathrm{trd}(wz)$ is diagonal on $1,i,j,ij$ with entries $2,-2d,2,2d$, so
the dual lattice $\cO^{\sharp}=\{w\in B:\mathrm{trd}(w\cO)\subseteq\ZZ\}$ is
$\ZZ\tfrac12\oplus\ZZ\tfrac{i}{2d}\oplus\ZZ\tfrac{j}{2}\oplus\ZZ\tfrac{ij}{2d}$.
And a form $y(u,v)=\sum_{k,l}\mathrm{trd}(\bbar{u_{l}}\,w_{lk}\,v_{k})$ on
$\cO^{2}$ with $w_{lk}\in B$ is integral if and only if every $w_{lk}$ lies
in $\cO^{\sharp}$, because $\mathrm{trd}(\bbar{u}wv)=\mathrm{trd}(wv\bbar{u})$
and the products $v\bbar{u}$ with $u,v\in\cO$ span $\cO$.

(i) The Gram matrix of $\mathrm{trd}(\bbar{u}iv)$ on $1,i,j,ij$ pairs $1$
with $i$ and $j$ with $ij$, with entries $\mp2d$ and $\pm2d$ and zeros
elsewhere, so $E$ has elementary divisors $2d$ four times on $\Lambda$ and
$e=4!\cdot(2d)^{4}=384\,d^{4}$. The involution $\sigma(x)=d^{-1}x(i\cdot,i\cdot)$
of $\bigwedge^{2}V^{*}$ has $R$ as its $(-1)$-eigenspace, since
$x(iu,v)=x(u,iv)$ is equivalent to $x(iu,iv)=-d\,x(u,v)$. Over $\CC$,
$V_{\CC}=W\oplus\bbar{W}$ with $W,\bbar{W}$ the eigenspaces of $i$ for
$\pm\sqrt{-d}$, each of dimension $4$, and $\sigma$ is $-1$ on
$\bigwedge^{2}W^{*}\oplus\bigwedge^{2}\bbar{W}^{*}$ and $+1$ on
$W^{*}\otimes\bbar{W}^{*}$; so $\dim R=6+6=12$. For $x\in R$ the map
$\phi_{x}$ is $E$-self-adjoint, as $x$ and $E$ are alternating, and
anticommutes with $i$, as $i$ is $E$-skew-adjoint by
\Cref{thm:quatweil}(i): $E(\phi_{x}iu,v)=x(u,iv)=E(\phi_{x}u,iv)=-E(i\phi_{x}u,v)$.
Hence $\mathrm{tr}\,\phi_{x}=-\mathrm{tr}(i\phi_{x}i^{-1})=-\mathrm{tr}\,\phi_{x}=0$,
and \Cref{lem:pfaffian}(i), with $g=4$, gives $t(x)=0$ and
$I(x)=-\int x^{2}\eta^{2}=\tfrac{1}{24}\mathrm{tr}(\phi_{x}^{2})e$.

For the signature regard $V_{\RR}$ as a complex vector space through
$i/\sqrt{d}$, with the hermitian form $H$ of \Cref{lem:hermform}, which has
signature $(2,2)$ by \Cref{thm:quatweil}(ii), and an $H$-orthogonal basis in
which $H=\mathrm{diag}(h_{1},\dots,h_{4})$. The maps $\phi_{x}$, $x\in R\otimes\RR$,
are the $E$-self-adjoint antilinear maps. Writing $\phi(v)=A\bbar{v}$,
self-adjointness for $E=d^{-1/2}\,\mathrm{Im}\,H$ says exactly that
$S=h\bbar{A}$ is antisymmetric, and then
$\mathrm{tr}_{\RR}(\phi^{2})=2\,\mathrm{Re}\,\mathrm{tr}(A\bbar{A})
=-4\sum_{k<l}|S_{kl}|^{2}/(h_{k}h_{l})$. Four of the six pairs $k<l$ have
$h_{k}h_{l}<0$ and two have $h_{k}h_{l}>0$, and each $S_{kl}$ is a complex
coordinate, so $I$ has signature $(8,4)$.

(ii) $\phi_{k}$ anticommutes with $i$ because $ji=-ij$ and the entries of
$C_{k}$ are rational. It is $E$-self-adjoint because $C_{k}$ is symmetric and
left multiplication by $j$ is self-adjoint for $\mathrm{trd}(\bbar{u}iv)$ on
each factor, by \Cref{thm:quatweil}(i); so $x_{k}\in R$. Since the entries of
$C_{k}$ commute with $j$, $\phi_{k}^{2}=(2d)^{-2}C_{k}^{2}j^{2}=(1+k^{2})/(4d^{2})$.
In the form of the first paragraph $x_{k}$ has
$w_{lk}=-(2d)^{-1}c_{kl}\,j\,i=(2d)^{-1}c_{kl}\,ij$, with $c_{kl}$ the entries of
$C_{k}$, which lie in $\cO^{\sharp}$; so $x_{k}$ is integral, and
$x_{k}\bigl((1,0),(ij,0)\bigr)=(2d)^{-1}\mathrm{trd}(ij\cdot ij)=1$, so it is
primitive.

(iii) Write $\phi=\phi_{k}$ and $\beta=\beta(k)$. The map $\phi$ is
$E$-self-adjoint, anticommutes with $i$, and $\phi^{2}=\beta>0$, so
$V_{\RR}=V_{+}\oplus V_{-}$ with $V_{\pm}$ the $\pm\sqrt{\beta}$ eigenspaces,
$E$-orthogonal to each other, and $i$ exchanges them. A complex structure $J$
on $V_{+}$ compatible with $E|_{V_{+}}$, extended to $V_{-}=iV_{+}$ by
$J(iw)=iJw$, commutes with $i$ and $\phi$ and is compatible with $E$, and $i$
has eigenvalues $\pm\sqrt{-d}$ each with multiplicity $2$ on $V^{1,0}$; so
these $J$ are points of $\cD_{2,2}$ at which $x_{k}$ is of type $(1,1)$, and
they form a copy of the Siegel space of $V_{+}$, of dimension $3$, in
agreement with \Cref{prop:lagrangian}. The algebra generated by $i$ and $\phi$
is $(-d,\beta)$, and $(-d,\beta)\cong(-d,1+k^{2})$ because the two differ by
the square $(2d)^{2}$. An automorphism $\gamma$ of $(\Lambda,E,i)$ sends
$\phi_{x}$ to $\gamma\phi_{x}\gamma^{-1}$ and so preserves
$\mathrm{tr}(\phi_{x}^{2})=8\beta$.

For the first half of (iv), let $S_{k}\subset\cD_{2,2}$ be the Siegel space
just described and $\mathcal{C}_{k}\subset\Sp(V_{\RR},E)$ the centraliser of
$i$ and $\phi$. An element of $\mathcal{C}_{k}$ preserves $V_{+}$ and is
determined by its restriction there, because it commutes with $i$, so
$\mathcal{C}_{k}\cong\Sp(V_{+})$, acting on $V_{\RR}=V_{+}\oplus iV_{+}$ as
$g\oplus igi^{-1}$. A real class is of type $(1,1)$ at $J$ exactly when it is
invariant under $J$. The points of $S_{k}$ form one orbit
$\{gJ_{0}g^{-1}:g\in\mathcal{C}_{k}\}$ of a fixed point $J_{0}$. Let $z$ be a
real $2$-form invariant under every point of $S_{k}$, and write
$\mathfrak{c}=\mathrm{Lie}(\mathcal{C}_{k})\cong\mathfrak{sp}(V_{+})$,
$\mathfrak{k}$ for the centraliser of $J_{0}$ in $\mathfrak{c}$ and
$\mathfrak{p}=[\mathfrak{c},J_{0}]$. For $X\in\mathfrak{c}$ the curve
$g_{t}=\exp(tX)$ gives $(g_{t}J_{0}g_{t}^{-1})^{*}z=z$ for all $t$, and the
derivative at $t=0$ says that $[X,J_{0}]$ annihilates $z$, so $\mathfrak{p}$
annihilates $z$. The annihilator of $z$ in $\mathfrak{c}$ is a Lie
subalgebra, so it contains $\mathfrak{p}+[\mathfrak{p},\mathfrak{p}]$. Now
$\mathfrak{c}=\mathfrak{k}\oplus\mathfrak{p}$ with
$[\mathfrak{k},\mathfrak{p}]\subseteq\mathfrak{p}$ and
$[\mathfrak{p},\mathfrak{p}]\subseteq\mathfrak{k}$, so
$\mathfrak{p}+[\mathfrak{p},\mathfrak{p}]$ is an ideal of $\mathfrak{c}$;
it is nonzero, and it is therefore all of $\mathfrak{c}$, because
$\mathfrak{c}\cong\mathfrak{sp}_{4}(\RR)$ is simple: its complexification
$\mathfrak{sp}_{4}(\CC)$ is simple, with irreducible root system $C_{2}$, and
an ideal of a real form complexifies to an ideal. Hence $z$ is annihilated by
$\mathfrak{c}$ and, since $\Sp(V_{+})$ is connected, invariant under all of
$\mathcal{C}_{k}$. So the classes of type $(1,1)$ at every point of $S_{k}$
are the invariants of $\Sp(V_{+})$ in $\bigwedge^{2}(V_{+}\otimes\RR^{2})
=\bigl(\bigwedge^{2}V_{+}\otimes S^{2}\RR^{2}\bigr)
\oplus\bigl(S^{2}V_{+}\otimes\bigwedge^{2}\RR^{2}\bigr)$, a space of dimension
$3$ spanned by $\eta$, $x_{k}$ and $\iota x_{k}$. Every other rational
class $z$ is of type $(1,1)$ only on a proper closed analytic subset of the
connected manifold $S_{k}$, and a very general point avoids these countably
many subsets.

For the lattice, $\iota x_{k}$ is the form of $\phi_{k}i=-i\phi_{k}$, so every
class in $\QQ\eta+\QQ x_{k}+\QQ\iota x_{k}$ is
$y=s\eta+E(\phi_{k}\lambda\,\cdot,\cdot)$ with $s\in\QQ$ and $\lambda\in K$
acting by left multiplication. In the form of the first paragraph, $y$ has
$w_{lk}=\delta_{lk}\,s\,i+(2d)^{-1}c_{kl}\bbar{\lambda}\,ij$. The first
summand lies in $K$ and the second in $Kj$, and $\cO^{\sharp}$ is the direct
sum of its parts in $K$ and in $Kj$, so $y$ is integral if and only if
$s\in(2d)^{-1}\ZZ$ and $(2d)^{-1}c_{kl}\bbar{\lambda}\,ij\in\cO^{\sharp}$
for all $k,l$. Writing $\bbar{\lambda}=p-q\,i$, one has
$\bbar{\lambda}\,ij=p\,ij+qd\,j$, and since $c_{11}=1$ and all $c_{kl}$ are
integers this holds exactly when $p,q\in\ZZ$, that is $\lambda\in\ZZ[i]$.
This is the stated description of $\NS$. Now $t(\eta)=1$ and $t$ vanishes
on $R$, so $I(s\eta+z)=s^{2}e+I(z)$ for $z\in R$, and
$(\phi_{k}\lambda)^{2}=\phi_{k}^{2}\bbar{\lambda}\lambda=N(\lambda)\beta(k)$,
so by (i)
\[
  I\bigl(s\eta+E(\phi_{k}\lambda\,\cdot,\cdot)\bigr)
  =s^{2}e+\tfrac{8}{24}N(\lambda)\beta(k)e
  =s^{2}e+32\,d^{2}(1+k^{2})N(\lambda).
\]
A class off $\QQ\eta$ has $\lambda\ne0$, hence $N(\lambda)\ge1$, and
$\lambda=1$, $s=0$ gives $x_{k}$; so $\mu_{k}=32\,d^{2}(1+k^{2})$. The final
assertion follows from \Cref{lem:nefnorm}(ii): an effective $D$ with
$[D]\notin\QQ\eta$ has
$(\int[D]\eta^{3})^{2}\ge e\,\mu_{k}/2=6144\,d^{6}(1+k^{2})$.
\end{proof}

Each locus of the pencil carries algebraic Weil classes, by the divisorial
criterion, since $x_{k}$ is nondegenerate, and on the $k$-th of them the
divisors that produce those classes have degree at least of order $|k|$.
This is \Cref{prop:divroute} seen on explicit quaternionic loci in one fixed
polarised lattice, and \Cref{fig:pencil} illustrates it.

\begin{figure}[t]
\centering
  \includegraphics[width=\textwidth]{fig_pencil}
\caption{\Cref{prop:pencil} at $d=3$. Left: the members $S_{k}$ of the
pencil, each a Siegel threefold inside the four-dimensional family, with the
member $k=0$, where $\beta=\tfrac{1}{36}$ is a square and a rational point
$J_{0}$ can be written down, drawn apart. Right: along the pencil the ratio
$\mu_{k}/\beta(k)$ is the constant $10368=\int\eta^{4}/3$, so $\mu_{k}$
grows like $k^{2}$ and the least degree of a divisor not proportional to
$\eta$ grows linearly in $k$. The Weil class is a flat class of the one fixed
lattice and does not change; the divisors that produce it must.}
\label{fig:pencil}
\end{figure}

\begin{remark}[The case $n=2$]\label{rem:ntwo}
Let $n=2$, so $\dim\cS=n^{2}=4$. For every $n$ the quaternionic locus of
\Cref{setup:quat} has dimension $n(n+1)/2$, by a Morita argument. The
quaternion algebra $B=(-d,b)$ contains $K$ and acts on $V$, and
$B\otimes\RR\cong M_{2}(\RR)$ because $b>0$, so Morita equivalence writes
$V_{\RR}$ as $W\otimes_{\RR}\RR^{2}$ with $W$ symplectic of dimension $2n$.
The complex structures commuting with $B$ and compatible with $E$ are then
the complex structures on $W$ compatible with its symplectic form, that is the
points of the Siegel
space $\mathfrak{H}_{n}$, of dimension $n(n+1)/2$, and the codimension in
$\cD_{n,n}$ is $n^{2}-n(n+1)/2=n(n-1)/2$. At
$n=2$ the quaternionic locus is therefore the symmetric space of
$\Sp(4,\RR)$ inside $\cD_{2,2}$, of dimension $3$, so it is a divisor. Its
$G(\QQ)$-translates are infinitely many
pairwise distinct irreducible closed subvarieties of $\cS$ of dimension $3$,
since a finite union of proper closed analytic subsets cannot be dense. If
infinitely many of them carry cycles with one and the same data
$(\underline{P},\underline{q})$, then the closed set
$\Sigma_{\underline{P},\underline{q}}$ has an irreducible component containing
infinitely many distinct irreducible threefolds, so that component has
dimension at least $4$ and therefore equals $\cS$. No density is needed at
$n=2$, only infinitude. Combined with the theorem of Moonen and Zarhin
\cite{MZ99} that every Hodge class on an abelian fourfold lies in the
subalgebra generated by divisor classes and Weil classes, the conclusion would
be the Hodge conjecture for all abelian fourfolds of Weil type, hence for all
abelian fourfolds. That conclusion is a theorem of Markman
\cite[Corollary~1.6.1]{Mar25a}, proved by a different route, so at $n=2$ the
argument above would give a second proof; its interest lies in $n\ge3$.
\end{remark}
